%%%%%%%%%%%%%%%%%%%%%%%%%%%%%%%%%%%%%%%%%%%%%%%%%%%%%%%%%%%%%%%%%%%%%%%%%
%                                                                       %
%                       A Quick Note on Alternatives                    %
%                                                                       %
%%%%%%%%%%%%%%%%%%%%%%%%%%%%%%%%%%%%%%%%%%%%%%%%%%%%%%%%%%%%%%%%%%%%%%%%%

% sometimes it is enough to just define a command,
% that just prints a formated heading and steps the counter with
% \refstepcounter{mysection}
% with a correct counter setup, this resets counters of childsections and
% \label and \ref work fine by calling \themysection


%%%%%%%%%%%%%%%%%%%%%%%%%%%%%%%%%%%%%%%%%%%%%%%%%%%%%%%%%%%%%%%%%%%%%%%%%
%                                                                       %
%                       Defining a New Type of Section                  %
%                                                                       %
%%%%%%%%%%%%%%%%%%%%%%%%%%%%%%%%%%%%%%%%%%%%%%%%%%%%%%%%%%%%%%%%%%%%%%%%%

% define a new counter for the new type of section
% the optional parentsection is a counter,
% that automatically resets the counter of mysection
\newcounter{mysection}[parentsection]


% representation of the counter
% this will be used in the custom section header and toc entry and references
\renewcommand\themysection{\normalfont\large{\bfseries MySection \arabic{mysection}.}}


% this command needs to be defined, as otherwise the argument is printed during startsection
\newcommand\teilnrmark[1]{}

\makeatletter
\newcommand{\mysection}[2]{ % The command to start the section. Does not need to match
    \@startsection          % The @startsection command is wraped. Parameters to @startsection are explained here
                            % Sometimes it is better to use currying by not consuming arguments and
                            % Siehe auch /usr/share/texmf-dist/tex/latex/base/latex.ltx, Zeile 17231
                            %
    {mysection}             % Counter, der für die section benutzt wird
                            %
    {1}                     % Tiefe der section
                            %
    {\z@}                   % Irgendwas mit indent
                            %
    {0em}                   % Platz davor?
                            %
    {-1em}                  % Platz danach, negative werte für horizontal anstatt vertikal
                            %
    {\emph}                 % Formatierung, wichtig wenn Referenzen verwendet werden
                            % und die Formatierung des counters in den Referenzen anders ist, als im Titel
                            %
    {#2}                    % Name der section. Kann auch gecurryed werden für
                            % optionale * und [toctitle] von startsection
                            %
    (insgesamt \punkte{#1}) % Text nach dem Gegeben Titel: hier Punktzahl der Aufgabe. Nicht möglich mit currying
    }%
\makeatother
